%    SeUnem - A free software platform for participatory decision-making and
%    collective event management.
%    Copyright (C) 2025  João Augusto Costa Branco Marado Torres
%    <torres.dev@disroot.org>
%
%    This file is part of SeUnem.
%
%    SeUnem is free software: you can redistribute it and/or modify it under
%    the terms of the GNU Affero General Public License as published by the
%    Free Software Foundation, either version 3 of the License, or (at your
%    option) any later version.
%
%    SeUnem is distributed in the hope that it will be useful, but WITHOUT ANY
%    WARRANTY; without even the implied warranty of MERCHANTABILITY or FITNESS
%    FOR A PARTICULAR PURPOSE.  See the GNU Affero General Public License for
%    more details.
%
%    You should have received a copy of the GNU Affero General Public License
%    along with SeUnem.  If not, see <https://www.gnu.org/licenses/>
\documentclass[12pt,a4paper,openright,oneside]{memoir}

\usepackage{iftex}
\ifXeTeX
  \usepackage{fontspec}
  \defaultfontfeatures{Ligatures=TeX}
  \setmainfont{EB Garamond}[
    Scale = 1.0
  ]
  % \setmonofont{Consolas}[
  %   Scale = 1.0
  % ]
\else
  \usepackage[T1]{fontenc}
  \usepackage[utf8]{inputenc}
  \usepackage{mathptmx}
\fi

\usepackage{polyglossia}
\setdefaultlanguage{portuguese}
\setotherlanguage{english}

\usepackage[a4paper,top=3cm,bottom=3cm,left=2cm,right=2cm]{geometry}

\OnehalfSpacing

\maxsecnumdepth{subsection}
\setcounter{secnumdepth}{3}
\setsecnumformat{\csname the#1\endcsname\quad}

\chapterstyle{section} % clean style
\renewcommand*{\chapnamefont}{\normalfont\Large\scshape}
\renewcommand*{\chaptitlefont}{\normalfont\Huge\bfseries}

% \usepackage{titlesec}
% \titleformat{\chapter}[display]
%   {\normalfont\huge\bfseries}{\chaptertitlename\ \thechapter}{20pt}{\Huge}
% \titleformat{\section}{\normalfont\Large\bfseries}{\thesection}{1em}{}
% \titleformat{\subsection}{\normalfont\large\bfseries}{\thesubsection}{1em}{}

% \usepackage{fancyhdr}
% \pagestyle{fancy}
% \fancyhf{}
% \fancyhead[L]{\fontsize{9}{11}\selectfont\leftmark}
% \fancyhead[R]{\fontsize{9}{11}\selectfont\rightmark}
% \fancyfoot[R]{\fontsize{8}{9}\selectfont\thepage}
% \renewcommand{\headrulewidth}{0.4pt}
% \renewcommand{\footrulewidth}{0pt}
% \nouppercaseheads
%
% \renewcommand{\chaptermark}[1]{\markboth{Capítulo \thechapter\ ---\ #1}{}}
% \renewcommand{\sectionmark}[1]{\markright{#1}}

\usepackage{caption}
\DeclareCaptionLabelFormat{ipbeja}{#1~#2}
\captionsetup{
  labelfont=bf,
  labelsep=none,
  format=plain,
  textfont=it,
  justification=justified,
  singlelinecheck=false,
  labelformat=ipbeja
}

\captionsetup[table]{position=top}

\usepackage{graphicx}
\usepackage{float}
\usepackage{booktabs}
\usepackage{longtable}

\usepackage{minted}
\setminted{
  linenos,
  breaklines,
  frame=lines,
  fontsize=\small
}

\usepackage[style=apa, backend=biber]{biblatex}
\addbibresource{bibliography.bib}
\usepackage{csquotes}

\newcommand{\Institute}{Instituto Politécnico de Beja}
\newcommand{\School}{Escola Superior de Tecnologia e Gestão}
\newcommand{\Degree}{Licenciatura em Engenharia Informática}
\newcommand{\Course}{Desenvolvimento de Aplicações Web}
\newcommand{\Title}{SeUnem}
\newcommand{\Subtitle}{Eventos Comunitários, Recursos Partilhados e Decisão Coletiva\mbox{ --- }Análise e Desenho}
\newcommand{\Author}{João Augusto Costa Branco Marado Torres}
\newcommand{\Advisor}{Dr. Luís Carlos da Silva Bruno}
\newcommand{\JuryMemberFirst}{Dr. Luís Filipe Nobre Horta Baptista Garcia}
\newcommand{\JuryMemberSecond}{Dr. João Paulo Barros}
\newcommand{\Date}{Beja, outubro de 2025}

\usepackage[hidelinks]{hyperref}
\usepackage{hyperxmp}
\hypersetup{
  pdfauthor={\Author},
  pdftitle={\Title},
  pdflicenseurl={https://creativecommons.org/licenses/by-sa/4.0/},
  pdfcopyright={© 2025 \Author --- CC BY-SA 4.0 for PDF, AGPL v3 for source},
}

\begin{document}

\thispagestyle{empty}

\begin{center}
    \includegraphics{Logotipo_IPBeja_horizontal-5/IPbeja_horizontal}

    \bigskip

    \textsc{\large \School}\\{\large \Degree}

    \bigskip

    \textsc{\large \Course}

    \vspace{4\baselineskip}

    \textsc{\Huge \Title}

    \smallskip

    {\Large \Subtitle}

    \bigskip

    {\large\bfseries \Author}

    \vfill

    \begin{center}
        \includegraphics[height=25mm,keepaspectratio]{Logotipos_ESTIG/estig}%
    \end{center}

    \vfill

    {\footnotesize \Date}
\end{center}

\cleardoublepage

\thispagestyle{empty}
\begin{center}
    \textsc{\large \Institute}

    \bigskip

    \textsc{\large \School}\\{\large \Degree}

    \bigskip

    \textsc{\large \Course}

    \vspace{4\baselineskip}

    \textsc{\Huge \Title}

    \smallskip

    {\Large \Subtitle}

    \bigskip

    {\large\bfseries \Author}

    \vspace{2\baselineskip}

    {\large Trabalho realizado no âmbito da unidade curricular de \Course}

    \vspace{2\baselineskip}

    \textsc{Orientação}

    \bigskip

    \Advisor

    \vfill

    {\footnotesize \Date}
\end{center}

\cleardoublepage

\thispagestyle{empty}
\begin{center}
  \textbf{Júri}

  \bigskip

  Responsável: \Advisor\\
  Vogal: \JuryMemberFirst\\
  Vogal: \JuryMemberSecond\\
\end{center}
\clearpage

\pagenumbering{roman}
% \fancyfoot[R]{\fontsize{8}{9}\selectfont\thepage}

\chapter*{Resumo}
\addcontentsline{toc}{chapter}{Resumo}
\noindent
...

\bigskip

\textbf{Palavras-chave:} Gestão comunitária; Eventos; Recursos partilhados;
Decisão coletiva; Participação democrática; Software MVC.

\chapter*{Abstract}
\addcontentsline{toc}{chapter}{Abstract}
\noindent
...

\bigskip

\textbf{Keywords:} Community management; Events; Shared resources; Collective
decision-making; Democratic participation; MVC software.

\chapter*{Dedicatória}
\addcontentsline{toc}{chapter}{Dedicatória}
\begin{flushright}
    ...
\end{flushright}

\chapter*{Agradecimentos}
\addcontentsline{toc}{chapter}{Agradecimentos}
...

\clearpage
\tableofcontents
\clearpage
\listoffigures
\clearpage
\listoftables
\clearpage

\chapter*{Lista de Abreviaturas e Siglas}
\addcontentsline{toc}{chapter}{Lista de Abreviaturas e Siglas}
\begin{description}
    \item[a11y] \textit{Accessibility}
    \item[API] \textit{Application Programming Interface}
    \item[DX] \textit{Developer Experience}
    \item[FLOSS] \textit{Free Libre and Open Source Software}
    \item[HTTP] \textit{Hypertext Transfer Protocol}
    \item[i18n] \textit{Internationalization}
    \item[MVC] \textit{Model-View-Controller}
    \item[POO] Programação Orientada a Objetos
    \item[REST] \textit{Representational State Transfer}
    \item[TDD] \textit{Test-Driven Development}
    \item[UML] \textit{Unified Modeling Language}
\end{description}

% \chapter*{Simbologia e Notação}
% \addcontentsline{toc}{chapter}{Simbologia e Notação}
% \begin{description}
%   \item[$x$] variável independente
%   \item[$y$] variável dependente
% \end{description}

\clearpage
\pagenumbering{arabic}
% \fancyfoot[R]{\fontsize{8}{9}\selectfont\thepage}

\chapter{Introdução}
\label{ch:intro}
% A introdução é um elemento essencial do trabalho que geralmente enquadra o
% objeto de estudo e apresenta os objetivos do trabalho, a metodologia
% utilizada e a estrutura do documento.

% Enquadramento geral do projeto, identificando-se o tema, os objetivos e a
% introdução ao sistema que se pretende desenvolver. Poderá usar o conteúdo
% submetido no moodle relativo à proposta de tema de projeto. Descrever
% sumariamente os capítulos que o relatório tem escritos posteriormente a este
% capítulo introdutório.

Este relatório é sobre a fase de análise e desenho do SeUnem, que por enquanto
é apenas uma plataforma de gestão de eventos democrática que também servirá
para eu ter uma ideia de como a tecnologia deve auxiliar a participação
democrática. Apenas auxiliar porque tenho a ideia que a participação
democrática é um problema social, e não vai ser com a tecnologia que esses
problemas vão ser resolvidos.\footcite{d3_vote} E como deve auxiliar porque só
pode existir uma forma correta de fazê-lo: usando \textit{Free Libre and Open
Source Software} (FLOSS) (o Software
Livre)\footcite{ansol_software_livre}\footcite{fsf_about_free_software}\footcite{gnu_free_software},
e serviços descentralizados (protocolos ao invés de
plataformas).\footcite{d3_platforms_protocols} Tem que ser possível o
\textit{self-host} em que o objetivo final seria ter algo
federado.\footcite{fedigov}

Recentemente ouvi falar dos projetos \autocite{decidim} e \autocite{voca}, que
apesar de nunca ter usado e não saber quase nada acerca deles, tenho a sensação
que muitas das ideias que queria aplicar no SeUnem foram também aplicadas lá, e
os temas são relativamente parecidos.

\section{Contexto}
% Apresente o contexto do estágio/projeto e a motivação.

Imagina que tu fases parte de uma comunidade. Era muita fixe se conseguissem
reunir-se para criar um evento e partilhar mais sobre as vossas paixões entre
vós, mas também talvez entre pessoas curiosas que viram um anúncio sobre o
vosso evento e quiseram aparecer. Eu nunca estive nessa situação para ser
sincero, mas não parece ser algo fácil de fazer. Uma comunidade é feita de
pessoas que têm algo em comum, mas não deixam de ser indivíduos com vidas
diferentes umas das outras. Eu gostava que local do evento fosse à frente da
minha casa, mas como eu vivo na Madeira e o resto da comunidade está espalhada
por Portugal continental, eles provavelmente iriam preferir que o evento fosse
algures entre a região Centro e a Península de Setúbal. Uns adoravam que as
atividades mais fixes fossem antes da hora de almoço, outros depois. É
complicado agradar todo o mundo, mas de certeza que com democracia direta, o
resultado para a maioria (senão todos) vai ser a satisfação. Como é que de
forma simples consegues convidar pessoas fora dessa bolha a fazer uma visita?
Mas não acaba por aí! Esses eventos vão requerer coisas, podem ser cadeiras e
mesas, ou hardware e software caso seja necessário disponibilizar ao vivo
streaming do evento para o bacano lá da ilha. Essa gestão de equipamentos não é
tão fácil, e é capaz até de ter custos um pouco elevados. Bom é que vós não
sois a única comunidade a existir em Portugal. Tem aquela comunidade meio
esquisita de [...] que no mês passado, fez o seu evento anual, e que por isso
já tem cadeiras e mesas disponíveis e que não estão em uso por mais uns 10
meses talvez, e que era fixe eles emprestarem esses produtos. Talvez só
precisem de metade das cadeiras que eles têm, bom. Talvez precisam das cadeiras
que eles têm e mais 20 (que lá em comunidade vão arranjar essas cadeiras que
falam), então arranjam as cadeiras que faltam e quem sabe para o ano essas
cadeiras novas vão fazer falta àqueles esquisitos.

Então o projeto vai ser um gestor de eventos para comunidades. Cria um evento
para uma comunidade, dá data, hora e local e uma pequena descrição. Depois
esses eventos serão listados. Se alguém quiser participar, diz que também vai,
e dá "\textit{up}"/"\textit{boost}" nesse evento (podendo ser usado para
apresentar os eventos mais populares). Quem vai e faz parte desse evento pode
dar feedback antes, durante, e depois do evento. Esse feedback pode fazer com
que data, hora ou local mudem. Se a mudança ocorrer, quem já ia participar tem
de ser notificado sobre a mudança para confirmar se ainda está interessado no
evento ou não (como padrão, ele não vai estar interessado até ele dizer que
sim). Como eventos normalmente necessitam de bens como cadeiras, mesas, entre
outras coisas, que podem muito bem ser reutilizadas de evento para evento,
também vai existir forma de "partilhar" esses bens. Se um evento acaba, mas
daqui a uma semana tem outro evento, as cadeiras e mesas podem ser reutilizadas
nesse outro evento, e é feita uma passagem de "\textit{ownership}" temporária
do bem. Quem participa no evento com bens em falta, vai ter palavra sobre quais
bens requisitar, da mesma forma que podem decidir mudanças de horário por
exemplo. Se existir bens em excesso, elas continuam no "banco", mas se faltar,
um dos organizadores compra o que falta e adiciona ao "banco" para puder ser
utilizado em outro evento futuro. Esta transferência de bens serve para fazer
\textit{analytics} também. Ver quais mercadorias tem maior disponibilidade e
quais são as mais necessárias. Deste modo, até alguém que não organiza os
eventos, pode doar mercadorias conforme as necessidades.

\section{Objetivos}
% Os objetivos gerais e específicos do trabalho.

Isto é um trabalho para a escola, e eu acredito que o ensino tem que ser livre,
outra razão para o projeto ser FLOSS também. Infelizmente não é isso que eu
vejo, preferem <<preparar-nos para trabalhar nas empresas>>. O software
proprietário é inerentemente contrário à educação e sabedoria, já que
\textbf{não deixa o estudante estudar} o programa, o que ele realmente faz e
como, modificar para as suas necessidades e a dos seus colegas e partilhar com
os mesmos a versão modificada. Nas aulas de Matemática e Física eu sempre tive
acesso à fórmulas e aos algoritmos, porquê que nas Ciências da Computação tem
que ser diferente? Afinal, \autocite{guimaraesPitagoras}! Eu vou fazer a minha
parte.

Também faz-me sentido que os dados e a informação sejam abertos. Para isso vou
ter em conta o \autocite{5stardata} e os \autocite{fair_principles}.

Temos que pensar na internacionalização (i18n) e ter acessibilidade (a11y) em
mente, pensar no offline, proteger a anonimidade dos utilizadores e encriptar
comunicações quando elas têm que existir.

Basicamente, a ideia não é existir diferentes tipos de utilizadores, todos os
utilizadores vão ter o mesmo “poder” em tudo. A pessoa que cria o evento, de
início, vai ter todo o poder sobre aquele evento. Mas depois de pessoas dizerem
que vão participar no evento, elas tem tanto quanto o criador. As decisões são
feitas por todos de forma democrática. Qualquer pessoa também poderá participar
num evento (ao menos que horários colidam, ou sei lá). Qualquer um também
consegue doar mercadorias para o banco. Mas pensando bem, era bom existir
pessoas que conseguem melhor gerir o banco de mercadorias. Porque após um
evento acabar, essas mercadorias não vão estar diretamente acessíveis, primeiro
vão precisar serem devolvidas ou banco. Aí era preciso um gajo para trocar o
status da mercadoria para “disponível”. Mas pensando mais um pouco, talvez não
tenha de ser um tipo de utilizador diferente. Pode existir um sistema de dizer
que alguém e o encarregado temporário da mercadoria, e ele tem a
responsabilidade de, ou por a mercadoria no banco e marcar o status da
mercadoria como disponível, ou entregar diretamente para o outro evento, e o
status continua indisponível. Nas duas situações ele deixa de ser o encarregado
da mercadoria. É de pensar.

Um dia talvez eu desenvolva o projeto além da ideia atual. Explorar, para além
da coordenação e mobilização comunitária e a troca de bens comuns e
distribuição circular, a situação do planeamento e orçamentação coletivos e do
conhecimento coletivo.

Promover a gestão coletiva, o combate a burocracia, e uma real democracia
direta.

\section*{Estrutura do documento}

O presente relatório encontra-se estruturado da seguinte forma:
\begin{itemize}
  \item O \textbf{Capítulo~\ref{ch:intro}} apresenta o enquadramento, a motivação, os objetivos e a metodologia do projeto;
  \item O \textbf{Capítulo~\ref{ch:analysis}} descreve a análise de requisitos, os atores e os casos de uso do sistema \textit{SeUnem};
  \item O \textbf{Capítulo~\ref{ch:design}} apresenta o desenho do sistema, incluindo diagramas UML e protótipos das interfaces;
  \item Por fim, o \textbf{Capítulo~\ref{ch:conclusion}} sintetiza as conclusões do trabalho e perspetivas de evolução futura.
\end{itemize}

\chapter{Desenvolvimento}
\label{ch:dev}

% O desenvolvimento consiste na parte nuclear do trabalho; descreve
% detalhadamente as ações realizadas, interpreta, ilustra, sistematiza e
% relaciona dados e resultados. O corpo do trabalho deve incluir, caso exista,
% de forma equilibrada e articulada com o texto, informação tabelar, imagens,
% gráficos, etc. O desenvolvimento subdivide-se em capítulos e, eventualmente,
% em subcapítulos.

\section{Análise do Sistema}
\label{ch:analysis}

% Em seguida descrevem-se o levantamento e a especificação de requisitos dos
% atores e dos casos de uso do sistema desenvolvidos durante esta fase de
% análise do sistema.

\subsection{Caracterização dos atores}

% Caracterização o perfil dos possíveis tipos de utilizadores, em várias
% vertentes (demográfica, físicas, social, literacia digital, entre outras) que
% possam ajudar a uma melhor compreensão dos utilizadores alvo do sistema.
% Especificar as personas como forma de melhor descrever esses perfis de
% utilizador.

Como talvez já deu para preceber, só existe um ator, coletivo, universal, e
humano. Sem importância para a sua raça, etnia, cultura e identidade. Sem olhar
para a carteira. No entanto, é de esperar que a pessoa tenha alguma prática com
a Web 2.0 para conseguir usar o software.

Alguém que gosta de conviver com outras pessoas, gosta de partilhar
conhecimento e que conhecimento seja partilhado consigo.

Talvez ainda dê para encaixar mais um ator que seria exclusivamente do sistema:
o banco de bens.

\subsection{Diagrama de Casos de Uso}

% Elaboração dos diagramas de casos de uso (5 no caso de um grupo de dois
% alunos e 2 no caso de um aluno) deste sistema seguido da especificação das
% respetivas tabelas de especificação de requisitos por cada caso de uso,
% usando a notação UML (Unified Modelling Language).

\begin{figure}[H]
  \centering
  \includegraphics[width=0.9\textwidth]{uml/analysis_n_design}
    \caption{Diagrama UML dos casos de uso: O actor \textit{User} listar
    eventos e agrupá-los, filtrá-los e ordená-los.}
  \label{fig:use_cases_list_event}
\end{figure}
\begin{figure}[H]
  \centering
  \includegraphics[width=0.9\textwidth]{uml/analysis_n_design_001}
    \caption{Diagrama UML dos casos de uso: O actor \textit{User} quando
    autenticado pode interagir com eventos seguindo-os e juntando-se a eles,
    criando-os, e participando democráticamente.}
    \label{fig:use_cases_interact_event}
\end{figure}
\begin{figure}[H]
  \centering
  \includegraphics[width=0.9\textwidth]{uml/analysis_n_design_002}
  \caption{Diagrama UML dos casos de uso: O actor \textit{User} pode e é
    encorajado a doar para o banco comun de bens.}
  \label{fig:use_cases_add_bank}
\end{figure}

\section{Desenho do Sistema}
\label{ch:design}

% Em seguida são descritos os modelos propostos de interfaces gráficos com o
% utilizador e do modelo de dados.

\subsection{Modelação da Base Dados}

% Descrição dos modelos E/R e relacional/físico da base de dados. O modelo E/R
% deve conter todas as entidades de dados e respetivas relações que descrevem
% conceptualmente o sistema. O diagrama relacional/físico descreve o a
% estrutura normalizada de tabelas e suas relações com a identificação dos seus
% atributos e respetivos tipos.

\subsubsection{Diagrama E/R \& Modelo Lógico Relacional}

\begin{figure}[H]
  \centering
  \includegraphics[width=0.9\textwidth]{uml/analysis_n_design_003}
  \caption{Modelo Entidade Relação.}
  \label{fig:er_diagram}
\end{figure}

\subsection{Modelação de Interfaces Gráficas com o Utilizador}

% Em seguida são descritos os modelos das interfaces gráficas desenvolvidas
% suportados em protótipo de baixa/média fidelidade das interfaces gráficas com
% o utilizador, compostos por um ou mais storyboard e os ecrãs (wireframes) da
% aplicação. O storyboard poderá ser desenvolvido preferencialmente em baixa
% fidelidade (papel e lápis) permite ter uma visão com baixo detalhe do esquema
% de apresentação das páginas e perceber genericamente a sequência de ações de
% navegação do utilizador no sistema. As propostas de wireframes
% (ecrãs/páginas) mostram as interfaces gráficas com o utilizador e devem ser
% agrupadas por caso de uso, mostrando o layout e o conteúdo de apresentação de
% cada página Web, com um nível de fidelidade alta, bem como, a forma como o
% utilizador interage com ela. As Wireframes podem ser desenvolvidas com
% ferramentas específicas para o efeito, como por exemplo, o Pencil, Balsamiq,
% InVision ou Figma.

Os protótipos de baixa e média fidelidade (\textit{wireframes}) foram criados
diretamente com \texttt{Svelte}, de forma a garantir uma representação fiel da
interação entre o utilizador e o sistema. Em vez de usar ferramentas externas
de prototipagem, optei por começar a desenvolver a estrutura das interfaces
reais, seguindo o princípio de que um protótipo executável é preferível a um
diagrama estático.

Todo o código-fonte correspondente encontra-se disponível no repositório
\autocite{repo}. Para reproduzir localmente os protótipos, podem seguir-se os
seguintes passos \ref{app:clone_wireframes}.

Após estes comandos, as interfaces podem ser visualizadas localmente através do
servidor de desenvolvimento. A versão deste commit (4f0d1c9) contém a estrutura
de páginas que serviu de base para a análise e desenho do sistema.

A opção por utilizar \texttt{Svelte} e \texttt{Inertia.js} (em vez de uma
ferramenta de prototipagem tradicional) permitiu uma abordagem mais
\textit{pragmática}, poupando tempo e assegurando a consistência entre o
protótipo e o código final da aplicação.

% \subsubsection{\textit{Storyboard(s)}}
% \subsubsection{Interfaces do caso de uso “1”}
% \subsubsection{Interfaces do caso de uso “2”}
% \subsubsection{Interfaces do caso de uso “3”}

\section{Metodologia}

Como conta para nota, tenho requisitos a cumprir, e por isso, não posso fazer
isto de qualquer forma. Vou ter que usar Programação Orientada a Objetos (POO)
para criar uma RESTful API e um sistema MVC (se não sabes de o que é que estou
a falar, vou tentar explicar no próximo parágrafo). Tive que escolher entre a
plataforma \autocite{laravel} e a \autocite{aspnet_core}. Apesar de hoje em dia
ASP.NET Core ser de código aberto, e eu sei que ia ser possível eu conseguir
fazer o trabalho no meu Arch Linux \footcite{tsoding}, continua a ser do
oligarca digital que é a Microsoft. Estarás mais sujeito a depender de algum
software ou serviço do ecossistema da Microsoft que não é nem código aberto e
muito menos FLOSS. Aliás, o que está no enunciado do projeto e o que dá a
perceber é que se a implementação do projeto for com ASP.NET, eu teria que
usar\mbox{ --- }vou ter que censurar\mbox{ --- }M\texttt{***}soft SQL
S\texttt{***}ver. Consegui fazer a cadeira de Bases de Dados 1 sem ter que
mexer nesse pedaço de software maligno e escapei de fazê-lo em Bases de Dados 2
porque fiz intercâmbio numa escola que não me obrigou a usar esse software. Mas
caraças, na cadeira de Sistemas de Informação também usam (junto com o
Microsoft Excel e outro que se chama Microsoft Power BI). Eu não tenho Windows
(e ainda bem) no meu portátil por isso só consigo fazer trabalhos que usem isso
nos PCs da escola. Nós não queremos Microsoft nas nossas vidas. Não é como se
ao escolher Laravel eu estivesse em melhores mãos porque a recomendação é MySQL
da Oracle, outro monopólio digital imperialista. MySQL tem a versão
\textit{Community} que é FLOSS, e a versão \textit{Enterprise} que é
proprietária. A chance de eu ter que usar as extensões proprietárias é mínima,
mas esse modelo de múltiplo licenciamento não quero suportar. Pessoas trabalham
de graça ao contribuir para a versão FLOSS e depois chega a Oracle que explora
esse trabalho em nome do capital. Para quem depende dessas extensões
proprietárias, agora estás preso\mbox{ --- }\textit{vendor lock-in}. Mas bem...
vou usar a \textit{fork} \autocite{mariadb_foundation}. Eu também já usei
\autocite{php} antes. É uma experiência agradável mesmo sem a tipagem estática,
e oiço falar de como Laravel tem um ótimo ecossistema e oferece uma ótima
\textit{developer experience} (DX), como \autocite{inertiajs}, que vai
permitir-me usar \autocite{svelte} para a parte das Views do sistema. Tem que
existir autenticação e um sistema de data analytics.

POO é um paradigma onde tu tens objetos, instâncias de classes. Uma classe é
como se fosse uma receita para criar as instâncias, que propriedades e
comportamentos elas terão. Depois cada instância do objeto é dona do seu estado
que não depende do estado dos outros objetos da mesma classe. Eu e tu
(assumindo que também és um Homo Sapiens Sapiens) ambos somos instâncias da
mesma "classe" mas com propriedades diferentes. Algumas dessas propriedades até
podem ser outros objetos, até de classes diferentes. Suponho que tu tenhas um
smartphone --- eu sei que tenho um. Tem a possibilidade de até termos o mesmo
modelo de smartphone, mas as propriedades individuais de ambos os smartphones
com certeza já são diferentes. Outra cena fixe é que objetos conseguem mandar
mensagens uns aos outros. Se eu carrego no botão para aumentar o volume do meu
telemóvel ele há de reagir de alguma forma (aumentando o volume espero eu). No
teu telemóvel deve acontecer o mesmo, mas não significa que aconteça da mesma
forma. Cada objeto instância da mesma classe, por ter sido usada a mesma
receita, oferecem os mesmos comportamentos, mas as propriedades por terem
valores diferentes, podem afetar a execução do comportamento. No final, é
esperado o mesmo resultado. Pensa assim, eu e tu provavelmente se nos mandassem
vestir as mesmas meias e sapatos, talvez eu comece na esquerda e tu na direita,
depois eu meto primeiro todas as meias e finalmente calço os 2 sapatos, e tu
calças o sapato logo depois de cada meia. Então essas mensagens podem ser
enviadas para qualquer tipo de objeto, instância da sua classe ou não, até para
o próprio objeto. As mensagens também podem ser enviadas e recebidas de forma
paralela.

API é uma interface que permite comunicação entre várias máquinas e/ou vários
programas de software. Essa interface oferece um serviço de um programa de
software de forma padronizada. O exemplo do botão do telemóvel, esse botão é
como uma interface.

REST é uma lista de princípios de estilo arquitetónico para software. Esses
princípios são na realidade restrições que limitam como um sistema de
\textit{hypermedia} distribuido deve agir. Podes estar agora a pensar que criar
restrições só deve complicar as coisas. Mas a verdade é que estas restrições
criam um sistema sem estado (\textit{stateless}) e fiável. Interfaces
uniformes, fácil separação de software em componentes menores, mais fácies de
manter. Se o teu sistema não quebra nenhuma das restrições REST, então ele é
\textit{RESTful}. Um dos pais do Hypertext Transfer Protocol (HTTP) foi quem
surgiu com a ideia do REST, na sua dissertação de maestrado
\autocite{fielding2000rest}.

Uma API \textit{RESTful} é uma API para a Web que é \textit{RESTful}. Existem
muitas formas de criar uma API \textit{RESTful}. O cliente pede um recurso de
\textit{hypermedia} e recebe uma representação do estado do recurso que ele
quis. Quando a API \textit{RESTful} é desenvolvida em cima do protocolo HTTP, o
método e o cabeçalho do \textit{request} pode servir como formar de dar
informação sobre tu e o teu \textit{request}, lembrando que a comunicação
servidor--cliente é proibída de guardar estado entre \textit{requests}.

MVC é uma padronização de arquitetura de software usando o paradigma de POO
falado em cima, que separa partes do código em 3 partes, cada uma com as suas
responsabilidades: o modelo, a vista, e o controlador. Um modelo representa
algum componente ou parte do sistema. As vistas são representações visuais dos
modelos. O controlador recebe mensagens do utilizador e manda-as tanto para os
modelos que possivelmente mandam mensagens para as vistas o que pode modificar
a informação que o utilizador vê (manda uma mensagem ao utilizador).

Nesta etapa eu vou caracterizar os atores do sistema e os casos de uso com que
eles interagem, usando \textit{Unified Modeling Language} (UML). Para isso,
usei o meu fiel \autocite{plantuml} com o documento PDF da última especificação
do UML\footcite{uml251} aberta.

Também querem um protótipo de baixa / média fidelidade das interfaces. Para
isso, o mais inteligente é só eu já começar a escrever o meu HTML, ou neste
caso, componentes do Svelte. Poupas trabalho para a próxima etapa.

Na minha opinião, quando tu não sabes como vai ser o teu projeto, mais vale
começar logo a programar, e não perder tempo com diagramas e coisas que, depois
que começas a programas, vais querer trocar. É uma perda de tempo. Tipo
\textit{Test-driven development} (TDD).

Para quê que tu queres <<uma visão com baixo detalhe do esquema de apresentação
das páginas>> quando podes ter algo feito com HTML 5 semântico e que já te
mostra <<a forma como o utilizador interage>> com as páginas Web e âncoras?

Matar dois coelhos de uma só cacheirada.

Depois gravo a tela e mando os vídeos como anexo junto ao relatório.

O professor também pede o modelo ER que eu nunca sei como fazer porque cada
professor faz de forma diferente e não existe uma padronização para eu seguir.
O PlantUML vai ajudar-me outra vez de qualquer forma.

O interessante deste projeto é que só existe um ator. Mas a interação entre
vários atores cria uma certa complexidade, especialmente na participação
democrática, que trás desafios já que todos os atores têm o mesmo poder.

O que também é de pensar é a questão da autenticação, porque dependendo de como
ela é feita, pode acabar com o sistema de democracia direta.

\subsection{Ferramentas}
\label{sec:ferramentas}

Para o desenvolvimento do projeto \textit{SeUnem}, foram utilizadas diversas
ferramentas livres e de código aberto, tanto para a implementação do sistema
como para a documentação e modelação. Entre as principais, destacam-se:

\begin{itemize}
  \item \textbf{Git} — Sistema de controlo de versões utilizado para gerir o
  histórico de desenvolvimento e facilitar o trabalho colaborativo.
  \item \textbf{Make} — Ferramenta de automatização utilizada para compilar e
  gerar a documentação técnica (diagramas e relatórios).
  \item \textbf{PlantUML} — Utilizada para a modelação dos diagramas UML,
  incluindo casos de uso e entidades-relacionamentos.
  \item \textbf{LaTeX} — Usado na redação e formatação do relatório técnico,
  com o auxílio do pacote \texttt{minted} para a inclusão de código.
  \item \textbf{PHP} — Linguagem de programação principal do \textit{backend}.
  \item \textbf{Laravel} — Framework PHP utilizada para estruturar a aplicação
  no lado do servidor, gerir as rotas e aceder à base de dados.
  \item \textbf{Composer} — Gestor de dependências do ecossistema PHP.
  \item \textbf{Node.js} e \textbf{pnpm} — Usados para gerir as dependências
  do \textit{frontend}.
  \item \textbf{Svelte} — Framework \textit{frontend} escolhida pela sua
  simplicidade e eficiência, integrada através de \texttt{Inertia.js}.
  \item \textbf{Vite} — Empregado como ferramenta de construção (\textit{build
  tool}) e servidor de desenvolvimento.
  \item \textbf{ESLint} e \textbf{Prettier} — Para assegurar a qualidade e
  consistência do código no lado do cliente.
  \item \textbf{MariaDB} — Base de dados leve usada durante o desenvolvimento e
  prototipagem inicial.
\end{itemize}

\subsection{Arquitetura do Sistema}
\label{sec:arquitetura}

A aplicação \textit{SeUnem} adota uma arquitetura de três camadas, composta por:

\begin{itemize}
  \item \textbf{Camada de Apresentação (Frontend)} — Implementada em
      \texttt{Svelte}, comunica com o servidor através do \texttt{Inertia.js},
        permitindo renderização dinâmica sem necessidade de uma API REST
        completa. Esta camada é responsável pela interação com o utilizador e
        pela exibição das informações de eventos, decisões e bens comunitários.

  \item \textbf{Camada Lógica (Backend)} — Desenvolvida em \texttt{Laravel},
      contém a lógica de negócio, autenticação, gestão de permissões e
        processos democráticos de decisão. O backend também gere a comunicação
        entre as outras camadas e implementa os princípios de autogestão e
        participação coletiva.

  \item \textbf{Camada de Dados (Base de Dados)} — \texttt{MariaDB}, é
      responsável pelo armazenamento persistente de entidades como
        utilizadores, eventos, decisões, votos e bens.
\end{itemize}

A comunicação entre o \textit{frontend} e o \textit{backend} ocorre através de
requisições HTTP geridas pelo \texttt{Inertia.js}, evitando a duplicação de
lógica e garantindo uma experiência fluida de aplicação de página única
(\textit{SPA — Single Page Application}). O sistema foi concebido de forma
modular, permitindo a extensão futura para novos tipos de entidades ou modos
de participação.

% \section{Resultados}
% Apresente resultados, gráficos e tabelas.

% \begin{table}[ht]
%   \centering
%   \caption{Número de visitantes do exemplo (tabela de exemplo).}
%   \label{tab:visitors}
%   \begin{tabular}{lrrr}
%     \toprule
%     Nome & Julho & Agosto & Média \\
%     \midrule
%     Castelo & 50 & 50 & 50 \\
%     Museu & 11 & 13 & 13.3 \\
%     \bottomrule
%   \end{tabular}
%   \vspace{0.3em}
%   \caption*{\textit{Nota.} Fonte: dados de exemplo.}
% \end{table}

% CHAPTER 3 - CONCLUSION
\chapter{Conclusão}
\label{ch:conclusion}
% Síntese dos trabalhos feitos, limitações e perspetivas futuras.

% A conclusão apresenta a síntese do trabalho realizado, podendo remeter para
% fases de estudo a desenvolver posteriormente. Corresponde, aproximadamente, a
% um décimo das páginas do corpo do trabalho. O seu cabeçalho é numerado.

% Resumo dos principais resultados obtidos nas tarefas desenvolvidas nesta fase
% do projeto, bem como, antecipando a implementação/programação do sistema por
% parte de cada autor do trabalho (em que tecnologia será desenvolvido cada
% grupo de casos de uso).

O desenvolvimento do projeto \textit{SeUnem} permitiu consolidar conhecimentos
adquiridos ao longo da licenciatura, nomeadamente no âmbito da análise,
desenho e implementação de sistemas Web baseados em paradigmas modernos de
engenharia de software. Esta fase de análise e desenho focou-se sobretudo em
definir a estrutura conceitual da aplicação, os atores envolvidos, os casos de
uso principais e o modelo de dados subjacente.

O sistema proposto tem como propósito promover a participação democrática,
autogestão e partilha comunitária através de software livre e descentralizado.
A escolha de tecnologias e ferramentas FLOSS garantiu coerência com a visão
ideológica do projeto — de que a tecnologia deve servir o coletivo e não o
contrário. A arquitetura em três camadas (frontend, backend e base de dados)
possibilita a evolução modular da plataforma, assegurando escalabilidade e
manutenção a longo prazo.

Entre as limitações encontradas, destaca-se o tempo reduzido para a
implementação efetiva das funcionalidades democráticas mais complexas (como o
processo deliberativo e de votação distribuída). Também se reconhece a
necessidade de aprofundar mecanismos de autenticação e privacidade que
preservem o anonimato sem comprometer a integridade da participação.

De forma geral, os objetivos desta etapa foram atingidos: elaboraram-se os
diagramas de casos de uso e E/R, a modelação da arquitetura, e os protótipos de
baixa/média fidelidade, demonstrando a viabilidade técnica e conceptual do
\textit{SeUnem}.

\section{Trabalho Futuro}

Numa fase futura, pretende-se evoluir o projeto no sentido de:

\begin{itemize}
  \item Implementar a totalidade da aplicação funcional baseada nos modelos
      definidos nesta fase;
  \item Desenvolver mecanismos de deliberação e votação descentralizada, com
      foco em transparência e auditabilidade pública;
  \item Melhorar a acessibilidade (\textit{a11y}) e a internacionalização
      (\textit{i18n}), tornando a plataforma inclusiva e global;
  \item Integrar protocolos federados para comunicação entre instâncias do
      \textit{SeUnem}, rumo a uma rede de comunidades autogeridas;
  \item Criar um módulo de \textit{analytics} participativo para visualizar a
      utilização dos bens e recursos comunitários;
  \item Estudar novas formas de autenticação descentralizada e mecanismos de
      reputação não-hierárquica.
\end{itemize}

Assim, o \textit{SeUnem} continuará a ser um espaço experimental e
tecnopolítico, orientado para repensar as práticas digitais de organização
coletiva e participação democrática, sempre com base em software livre e
valores de autonomia e partilha.

%%%



\clearpage
\printbibliography[title={Referências Bibliográficas}]

\clearpage
\chapter*{Licença}
\addcontentsline{toc}{chapter}{Licença}
\noindent
Este documento está licenciado sob uma
\href{https://creativecommons.org/licenses/by-sa/4.0/}{Licença Creative Commons Atribuição–Partilha nos Mesmos Termos 4.0 Internacional (CC BY-SA 4.0)}.

\vspace{0.5cm}
O código fonte (ficheiros \texttt{.tex}, \texttt{.bib}, \texttt{Makefile}, etc.) utilizado para produzir este relatório está licenciado sob a
\href{https://www.gnu.org/licenses/agpl-3.0.html}{GNU Affero General Public License v3.0 (AGPL v3)}.

\clearpage
\appendix
% \chapter{Apêndice A: Questionário}
% \label{app:questionario}
% Conteúdo do apêndice (ex.: questionários, listagens de dados).

\chapter{Apêndice: Linha de comandos}
\label{app:clone_wireframes}
\begin{minted}[linenos, frame=lines, fontsize=\small]{sh}
# Clonar o repositório
git clone https://github.com/torres-engineer/SeUnem.git seunem
cd seunem

# Alterar para o commit específico dos wireframes
git checkout 4f0d1c9

# Instalar as dependências do frontend e construir os ficheiros
pnpm install
pnpm build

# Iniciar o ambiente de desenvolvimento do Laravel
composer run dev
\end{minted}

\backmatter
% \chapter*{Anexos}
% \addcontentsline{toc}{chapter}{Anexos}
% \section*{Anexo A: Documento Externo}
% Texto de anexo.

\end{document}

